\documentclass[12pt]{article}
\usepackage[margin=0.65in]{geometry}

\newcommand{\duedate}{02/10/2026}
\newcommand{\assignment}{Recurrent Memory Transformer}

\newcommand{\name}{Aksel Kretsinger-Walters, adk2164}
\newcommand{\email}{adk2164@columbia.edu}

\makeatletter
\def\input@path{{../}{../../}{../../../}}
\makeatother
\input{pset_template_CLMM.tex} %% DO NOT CHANGE THIS LINE

\begin{document}
\psetheader %% DO NOT CHANGE THIS LINE

\section*{Recurrent Memory Transformer}

\paragraph{Motivation.}
Standard Transformers mix local token representations with global context and scale quadratically in sequence
length, while Transformer-XL extends context via cached hidden states but entangles token and memory
representations and restricts gradient flow across segments.

\paragraph{Core Idea.}
The Recurrent Memory Transformer (RMT) augments a vanilla Transformer with dedicated memory tokens that are
prepended for reading and appended for writing, allowing explicit storage and recurrent propagation of
compressed sequence information across segments.

\paragraph{Memory Mechanism.}
At each segment, memory tokens are processed jointly with input tokens through all Transformer layers, and
the output memory tokens are passed forward as the memory state for the next segment, enabling segment-level
recurrence without modifying internal Transformer layers.

\paragraph{Architectural Distinction.}
Unlike Transformer-XL, which caches layer-wise hidden states, RMT uses a fixed number of explicit memory
vectors, reducing representation mixing and enabling deeper memory processing as the same memory tokens
traverse the network repeatedly.

\paragraph{Backpropagation Through Time.}
RMT allows gradients to flow across segments via truncated backpropagation through time, improving long-range
credit assignment compared to stop-gradient mechanisms in Transformer-XL.

\paragraph{Empirical Findings.}
On synthetic sequence tasks such as Copy and Reverse, RMT maintains accuracy across long segment chains where
Transformer-XL degrades, and on WikiText-103 it achieves competitive perplexity with significantly smaller memory sizes.

\paragraph{Interpretation.}
RMT can be viewed as a recurrent compression mechanism that separates transient token processing from
persistent storage, enabling efficient long-context modeling while preserving compatibility with pretrained
Transformer architectures.

\paragraph{Key Takeaway.}
By introducing explicit, recurrent memory tokens without altering Transformer internals, RMT provides a
simple yet powerful framework for long-sequence modeling based on structured memory rather than hidden-state caching.

\end{document}
